\section{The Body That Has to Come Back}\label{sec:body}

The redundancy objection has an unexamined premise. It assumes that the two judgments have the same defendant. If they do, then the second is a rehearing, and rehearing a case whose verdict is already in force and already executed really is empty ceremony. But they do not have the same defendant, and the reason they do not is the one part of the last day that no one has ever proposed to allegorize.

\subsection{The resurrection is the precondition, not an accompaniment}

The Catechism's treatment of the last judgment opens with a sequence, not a scene: ``The resurrection of all the dead, `of both the just and the unjust,' will precede the Last Judgment.''\endnote{\emph{Catechism of the Catholic Church} 1038, official English, read 25 July 2026 at the registered archived page; the internal quotation is Acts 24:15, and the paragraph continues with John 5:28--29 and Matthew 25:31--32, 46. The word ``precede'' is the Catechism's own.} The resurrection is not scenery for the judgment and not a consequence of it. It comes first, and the judgment presupposes it.

Scripture puts the two together in exactly this order. John's second hour is a resurrection before it is a judgment: all that are in the graves hear the voice, and \emph{then} they come forth, some to the resurrection of life and some to the resurrection of judgment. Paul's sequence in 1~Corinthians~15 is a sequence of risings---Christ the firstfruits, then they that are of Christ at his coming, then the end. The Apocalypse has the sea and death and hell give up their dead \emph{before} the books are read. And Lateran~IV, defining the last day, defines it as a rising: all shall rise \emph{cum suis propriis corporibus, quae nunc gestant}, that they may receive according to their works. The Catechism quotes that clause in its own English---``all of them will rise again with their own bodies which they now bear''---and sets beside it Paul's promise that Christ ``will change our lowly body to be like his glorious body.''\endnote{\emph{Catechism of the Catholic Church} 999, official English, read 25 July 2026 at the registered archived page; the paragraph quotes Luke 24:39, Lateran~IV (DS 801), Philippians 3:21, and 1~Corinthians 15:35--53. Paragraph~997 gives the anthropology the argument of this section uses: ``In death, the separation of the soul from the body, the human body decays and the soul goes to meet God, while awaiting its reunion with its glorified body.'' Paragraph~1001 answers the question ``When?'' with ``Definitively `at the last day,' `at the end of the world.'\,''}

The two adjectives in the definition are doing work: the bodies are \emph{their own}, and they are the ones they \emph{now bear}. Whatever transformation the risen body undergoes---and Paul's language of a spiritual body, and the Catechism's frank ``this `how' exceeds our imagination and understanding,'' both insist that it undergoes one---the identity of the risen man with the man who died is not negotiable. The resurrection is not the production of a new subject.

\subsection{Who stands at each judgment}

Now put the two judgments side by side and ask what is present at each.

At death, what receives its sentence is a soul. The Catechism says so with unusual precision: ``Each man receives his eternal retribution in his \emph{immortal soul} at the very moment of his death.'' \emph{Benedictus Deus} says the same---it defines what happens to \emph{animae}, souls, and it specifies that this occurs \emph{ante resumptionem suorum corporum}. The particular judgment is real and its effect is real, and its subject is one part of a man.

At the last day, what stands is a man---soul and the body he now bears, reunited. That is not the same defendant in the sense that matters for the objection. It is the same person, of course; it is not the same standing.

Aquinas draws exactly this consequence, and states it as the reply to the sharpest form of the objection. Grant that merit and demerit belong to the body only as the soul's instrument, and that reward or punishment is therefore due to the body only through the soul; does not that make a bodily judgment superfluous? His answer:

\begin{quote}
Although the reward or punishment of the body depends upon the reward or punishment of the soul, nevertheless, since the soul is changeable only accidentally, on account of the body, once it is separated from the body it enters into an unchangeable condition, and receives its judgment. But the body remains subject to change down to the close of time: and therefore it must receive its reward or punishment then, in the last Judgment.\endnote{\emph{Summa theologiae} III, q.~59, a.~5, ad~3, English Dominican translation, witness as in Section~\ref{sec:twice}. Latin at the registered Corpus Thomisticum page: \emph{Sed corpus remanet mutabilitati subiectum usque ad finem temporis. Et ideo oportet quod tunc recipiat suum praemium vel poenam in finali iudicio.} The argument turns on Aquinas's general principle, stated in the corpus of the same article, that perfect judgment cannot be passed on a changeable subject before its consummation.}
\end{quote}

That is the pivot of the whole question and it deserves to be stated in plainer words. A soul, once separated, cannot change; therefore it can be judged perfectly at once, and is. A body can change, and goes on changing---it is buried or left unburied, honoured or desecrated, and finally falls to dust---until time itself stops. A subject that is still changing cannot be judged perfectly. Therefore the part of a man that is still in motion is judged when its motion ends, which is when everything's motion ends.

\subsection{What the general judgment adds}

If the general judgment does something rather than merely announcing something, what does it add? The tradition's answer, and it is a common theological answer rather than a defined one, is: an increase.

\begin{quote}
Even now, as regards the particular sentence on each individual, the judgment does not at once take full effect since even the good will receive an increase of reward after the judgment, both from the added glory of the body and from the completion of the number of the saints. The wicked also will receive an increase of torment from the added punishment of the body and from the completion of the number of damned to be punished.\endnote{\emph{Supplementum} q.~88, a.~1, ad~2, English Dominican translation read 25 July 2026 at the registered New Advent page. The corresponding Latin in Aquinas's own \emph{Scriptum super Sententiis} IV, d.~47, q.~1, a.~1, qc.~1, ad~2, read the same day at the registered Corpus Thomisticum page, runs: \emph{quia etiam boni amplius post judicium praemiabuntur, tum ex gloria corporis adjuncta, tum ex numero sanctorum completo}. This is common theological teaching and not a defined proposition; nothing in Lateran~IV, Lyons~II, \emph{Benedictus Deus}, Florence, or Trent defines an increase of the essential reward or punishment at the general judgment, and the tradition holds the essential reward---the vision of God---to be already possessed. The increase asserted here is accidental. That distinction is load-bearing and is examined again in the synthesis block in Section~\ref{sec:synthesis}.}
\end{quote}

Two increases are named, and they are of different kinds. The first is the glory or punishment of the body, which is straightforward: an embodied person can be glorified in a way a separated soul cannot. The second is stranger and more interesting: the \emph{completion of the number}. The blessed are not complete while the roll is unfinished; each is happier for the others' happiness, and the happiness of all is not yet assembled. The remark about the damned is the same principle inverted and is stated with a bleakness the English does not soften: ``the more numerous those with whom they will burn, the more will they themselves burn.''

Whether one finds that persuasive, it is a genuine answer to the redundancy objection and not an evasion of it. It concedes that the essential reward is already given and locates the addition elsewhere. That concession is the strongest card the objector has, and Section~\ref{sec:trial} will hand it to him at full strength.

\subsection{The interval, and what may not be filled in}

Between a man's death and the resurrection there is an interval. What can be said about it is much less than what is usually said.

The medieval treatment devotes a question to the ``receptacles'' of souls and reasons about them with considerable care and considerable hedging: incorporeal things are not in place ``after a manner known and familiar to us,'' but only ``after a manner befitting spiritual substances, a manner that cannot be fully manifest to us.'' The souls that have a perfect share of the Godhead are in heaven and those deprived of it in a contrary place; and the movement is immediate, ``unless it be held back by some debt, for which its flight must needs be delayed until the soul is first of all cleansed.''\endnote{\emph{Supplementum} q.~69, a.~1, ad~1 and corpus, and a.~2, corpus, English Dominican translation read 25 July 2026 at the registered New Advent page. The image in a.~2 is of gravity: as a body is carried at once to its place by its weight unless obstructed, so the soul, ``the bonds of the flesh being broken,'' receives at once its reward or punishment unless obstructed. Article~2's \emph{sed contra} uses 2~Corinthians 5:1, Philippians 1:23 with Gregory's gloss, and Luke 16:22, and the corpus closes by calling the contrary opinion heretical on the authority of Gregory's \emph{Dialogues} and the \emph{De ecclesiasticis dogmatibus}. The status of this text is established in Section~\ref{sec:trial}.} The apparatus of limbos and abodes that occupies the rest of that question is school theology, some of it about questions the Church has since declined to settle, and none of it is used here.

What the Church has said in the modern period about this interval is short and mostly negative. The 1979 letter affirms that ``a spiritual element survives and subsists after death, an element endowed with consciousness and will, so that the `human self' subsists,'' and that the Church calls this element the soul and sees no valid reason to reject the word. It affirms that the resurrection refers to \emph{the whole person}. It affirms that the Lord's glorious manifestation is ``distinct and deferred with respect to the situation of people immediately after death.'' And it excludes ``every way of thinking or speaking that would render meaningless or unintelligible her prayers, her funeral rites and the religious acts offered for the dead''---which is a doctrinal argument from the liturgy, and a strong one: a Church that prays for the dead is committed to there being dead to pray for, in a condition in which prayer can reach them.\endnote{Sacred Congregation for the Doctrine of the Faith, \emph{Recentiores episcoporum Synodi}, 17~May 1979, points~2, 3, 4, and 5, Holy See English, read 25 July 2026 at the registered page. Point~4 adds that these prayers, rites, and acts ``are, in their substance, \emph{loci theologici}.'' The letter describes the situation it addresses as one in which ``the question is put of what happens between the death of the Christian and the general resurrection'' and in which discussions ``about the existence of the soul and the meaning of life after death'' had disturbed the faithful; it does not name any theologian, and neither does this article, which read no author's proposal at first hand and therefore attributes none.}

One point in the letter is easy to skim past and settles something. The sixth point excludes any explanation of man's destiny after death ``that would deprive the Assumption of the Virgin Mary of its unique meaning, namely the fact that the bodily glorification of the Virgin is an \emph{anticipation} of the glorification that is the destiny of all the other elect.'' If bodily glorification before the general resurrection is a unique privilege in one case, then in every other case it is genuinely deferred. The Assumption, which might look like an eschatological curiosity, is in fact a control on the whole doctrine: it is what makes the interval real.

And the letter is as firm as Trent would be about imagination:

\begin{quote}
When dealing with man's situation after death, one must especially beware of arbitrary imaginative representations: excess of this kind is a major cause of the difficulties that Christian faith often encounters. Respect must however be given to the images employed in the Scriptures. Their profound meaning must be discerned, while avoiding the risk of over-attenuating them, since this often empties of substance the realities designated by the images. Neither Scripture nor theology provides sufficient light for a proper picture of life after death.\endnote{Same letter, Holy See English, read 25 July 2026. The pastoral section adds the instruction that catechists must be ``careful not to allow childish or arbitrary images to be considered truths of faith.'' The double caution---against inflating the images and against emptying them---is the letter's own and is the standard this article has tried to observe.}
\end{quote}

Both halves of that are needed. The images may not be inflated into information, and they may not be deflated into feelings. What remains, when both errors are refused, is a doctrine with a small number of firm joints and a great deal of declared darkness between them. The next section takes the firmest joint---that the second judgment is not a second trial---and shows why.
